% ==========================================================================
%  Guia de Conteúdo — Prova 2 (2026)
%  Macroeconomia I (PPGEco-UFES, 2026)
%  Prof. Dr. Ricardo Ramalhete Moreira
%  Base: Lista 3 2026 ("lista da Prova 2") + CHAVE oficial de respostas
%        + slides das Aulas 10-13 (Romer caps. 8, 9, 11 e 12)
%  Complemento: provas/p2_simulado_2026.pdf (simulado no modelo da P1)
% ==========================================================================
\documentclass[12pt, a4paper]{article}

\usepackage[utf8]{inputenc}
\usepackage[T1]{fontenc}
\usepackage{lmodern}
\usepackage[brazil]{babel}
\usepackage{microtype}

\usepackage[top=2.5cm, bottom=2.5cm, left=3cm, right=3cm, headheight=15pt]{geometry}

\usepackage{amsmath, amssymb, amsthm}
\usepackage{mathtools}
\usepackage{bm}

\usepackage{booktabs}
\usepackage{array}
\usepackage{multicol}
\usepackage{xcolor}
\usepackage{hyperref}
\usepackage{enumitem}
\usepackage{tcolorbox}
\tcbuselibrary{breakable}
\usepackage{fancyhdr}
\usepackage{titlesec}

% --------------------------------------------------------------------------
% Cores
% --------------------------------------------------------------------------
\definecolor{cyanrbc}{HTML}{3b9dbd}
\definecolor{orangerbc}{HTML}{d97b39}
\definecolor{grayrbc}{HTML}{6b7280}
\definecolor{greenbox}{HTML}{166534}
\definecolor{greenbg}{HTML}{dcfce7}
\definecolor{boxbg}{HTML}{f5f0e8}
\definecolor{boxborder}{HTML}{3b9dbd}
\definecolor{provabg}{HTML}{fff7ed}
\definecolor{provaborder}{HTML}{d97b39}
\definecolor{correctbg}{HTML}{dcfce7}
\definecolor{correctborder}{HTML}{166534}

% --------------------------------------------------------------------------
% Caixas
% --------------------------------------------------------------------------
\newtcolorbox{resultbox}[1][]{standard, colback=boxbg, colframe=boxborder,
  boxrule=1pt, arc=3pt, left=6pt, right=6pt, top=4pt, bottom=4pt,
  fonttitle=\bfseries\small, breakable, #1}
\newtcolorbox{provabox}[1][]{standard, colback=provabg, colframe=provaborder,
  boxrule=1.2pt, arc=3pt, left=6pt, right=6pt, top=4pt, bottom=4pt,
  fonttitle=\bfseries\small, breakable, #1}
\newtcolorbox{correctbox}[1][]{standard, colback=correctbg, colframe=correctborder,
  boxrule=1pt, arc=3pt, left=6pt, right=6pt, top=4pt, bottom=4pt,
  fonttitle=\bfseries\small\color{correctborder}, breakable, #1}
\newtcolorbox{formulabox}[1][]{standard, colback=greenbg!30, colframe=greenbox!60,
  boxrule=0.8pt, arc=3pt, left=6pt, right=6pt, top=4pt, bottom=4pt,
  fonttitle=\bfseries\small, breakable, #1}

% --------------------------------------------------------------------------
% Comandos matemáticos
% --------------------------------------------------------------------------
\newcommand{\E}{\mathbb{E}}
\newcommand{\dif}{\mathrm{d}}
\newcommand{\VAR}{\mathrm{VAR}}

\setlength{\parskip}{0.5\baselineskip}
\setlength{\parindent}{0pt}

\titleformat{\section}{\large\bfseries\color{cyanrbc}}{Tópico \thesection.}{0.5em}{}[\vspace{-0.3em}\rule{\linewidth}{0.4pt}]
\titleformat{\subsection}{\normalsize\bfseries\color{grayrbc}}{}{0em}{}

\pagestyle{fancy}
\fancyhf{}
\lhead{\footnotesize\textcolor{grayrbc}{Guia de Conteúdo P2 2026 --- Macroeconomia I (PPGEco-UFES)}}
\rhead{\footnotesize\textcolor{grayrbc}{Romer caps.\ 8 $\cdot$ 9 $\cdot$ 11 $\cdot$ 12}}
\cfoot{\small\thepage}
\renewcommand{\headrulewidth}{0.4pt}

% ==========================================================================
\begin{document}
% ==========================================================================

% --------------------------------------------------------------------------
% Capa
% --------------------------------------------------------------------------
\begin{titlepage}
\begin{center}
\vspace*{1.2cm}
{\Large\bfseries Universidade Federal do Espírito Santo}\\[0.5em]
{\large Programa de Pós-Graduação em Economia (PPGEco-UFES)}\\[2em]

\rule{\linewidth}{1.2pt}\\[0.8em]
{\LARGE\bfseries\color{cyanrbc} Guia de Conteúdo --- Prova 2\\[0.4em]
\large Macroeconomia I --- 2026}\\[0.4em]
{\normalsize\color{grayrbc} Consumo $\cdot$ Investimento $\cdot$ Inflação e Política Monetária $\cdot$ Política Fiscal}\\[0.8em]
\rule{\linewidth}{1.2pt}

\vspace{1em}
{\large Estruturado pela Lista 3 (2026) e sua chave oficial de respostas,\\
integrado ao conteúdo dos slides das Aulas 10--13}

\vspace{1.5em}
\begin{tcolorbox}[standard, colback=provabg, colframe=provaborder, width=0.92\textwidth,
  arc=4pt, boxrule=1.2pt, left=8pt, right=8pt, top=6pt, bottom=6pt]
\raggedright
\textbf{\textcolor{provaborder}{Mapa do guia} (15 tópicos, seguindo as 8 questões da lista):}\\[0.4em]
\textbf{Parte I --- Consumo} (Aula 10; Romer cap.\ 8): Tópicos 1--4
\hfill $\leftarrow$ Lista Q1, Q2, Q3\\
\textbf{Parte II --- Investimento} (Aula 11; cap.\ 9): Tópicos 5--7
\hfill $\leftarrow$ Lista Q4, Q5\\
\textbf{Parte III --- Inflação e Pol.\ Monetária} (Aula 12; cap.\ 11): Tópicos 8--12
\hfill $\leftarrow$ Lista Q6\\
\textbf{Parte IV --- Déficits e Pol.\ Fiscal} (Aula 13; cap.\ 12): Tópicos 13--15
\hfill $\leftarrow$ Lista Q7, Q8
\end{tcolorbox}

\vspace{1.2em}
\begin{tcolorbox}[standard, colback=boxbg, colframe=boxborder, width=0.92\textwidth,
  arc=4pt, boxrule=1pt, left=8pt, right=8pt, top=6pt, bottom=6pt]
\raggedright
\textbf{Como este guia foi montado:} cada tópico apresenta (i) a teoria tal
como dada em aula, com as derivações; (ii) a \textbf{resposta oficial da
chave} da Lista 3, destacada em caixa verde --- é o padrão de resposta que o
professor considera correto; e (iii) uma caixa laranja \textbf{``Como cai na
prova''}, com os itens V/F que já caíram na P2 2021 e na P2 2024 sobre aquele
ponto. A prova será no modelo da P1 (julgar afirmativas, escolha única) ---
treine com o simulado em \texttt{provas/p2\_simulado\_2026.pdf}.
\end{tcolorbox}

\vfill
{\small Fontes: Lista 3 2026 + CHAVE oficial; slides das Aulas 10--13
(caps.\ 8, 9, 11 e 12 de Romer, 2012); P2 2021 (prova e respostas);
P2 2024 (chave oficial). Uso pessoal de estudo.}
\end{center}
\end{titlepage}

% ==========================================================================
% PARTE I — CONSUMO
% ==========================================================================
\begin{center}
{\LARGE\bfseries\color{orangerbc} Parte I --- Teoria do Consumo}\\[0.2em]
{\normalsize\color{grayrbc} Aula 10 $\cdot$ Romer cap.\ 8 $\cdot$ Lista 3: Q1, Q2 e Q3}\\[0.2em]
\rule{\linewidth}{0.8pt}
\end{center}

% ------------------------------------------------------------------
\section{Hipótese da Renda Permanente (Friedman, 1957)}
% ------------------------------------------------------------------

\subsection{Lista 3, Q1 (caput e item a) --- o problema e sua solução}

O consumidor vive $T$ períodos, com utilidade
$U = \sum_{t=1}^{T} u(C_t)$, $u'>0$, $u''<0$, taxa de juros nula e restrição
orçamentária $\sum_{t=1}^{T} C_t \leq A_0 + \sum_{t=1}^{T} Y_t$
($A_0$ = riqueza inicial). O Lagrangiano é
\[
L = \sum_{t=1}^{T} u(C_t)
  + \lambda\Bigl(A_0 + \sum_{t=1}^{T} Y_t - \sum_{t=1}^{T} C_t\Bigr).
\]
A condição de primeira ordem para cada $C_t$ é $u'(C_t) = \lambda$: como
$\lambda$ é o mesmo em todos os períodos e $u'$ é estritamente decrescente, o
consumo é \textbf{constante} no tempo ($C_1 = C_2 = \cdots = C_T$).
Combinando com a restrição ativa:

\begin{formulabox}[title={Função consumo da renda permanente}]
\[
C_t = \frac{1}{T}\Bigl(A_0 + \sum_{\tau=1}^{T} Y_\tau\Bigr)
\qquad\text{para todo } t.
\]
O termo entre parênteses são os \emph{recursos de vida inteira}; o consumo de
cada período é a fração $1/T$ deles (\textbf{suavização do consumo}). O lado
direito é a \textbf{renda permanente}; a diferença entre renda corrente e
renda permanente é a \textbf{renda transitória}.
\end{formulabox}

\begin{correctbox}[title={Chave oficial (Q1a): diferença para a hipótese keynesiana}]
\emph{``A principal diferença para com a teoria convencional keynesiana está no
fato de que esta associa de modo estável o consumo familiar à renda disponível
corrente, omitindo o papel das expectativas do consumidor quanto à trajetória
futura de sua renda. Já na teoria da renda permanente, o fluxo futuro esperado
de renda possui peso mais relevante, visto que qualquer incremento na renda
disponível corrente será ajustado ou diluído, tendo menor impacto à medida que
o $T$ de avaliação cresce.''}

Em números: um aumento \textbf{transitório} de $Y_1$ tem impacto marginal
$1/T$ sobre o consumo corrente (efeito decrescente em $T$); um aumento
\textbf{permanente} (replicado em $Y_2,\dots,Y_T$ na mesma magnitude) soma $T$
parcelas e eleva o consumo praticamente \textbf{um para um}. Por isso a
propensão marginal a consumir de curto prazo é \emph{menor} que a keynesiana
--- e um corte de impostos \emph{temporário} tem pouco efeito sobre o consumo.
\end{correctbox}

\subsection{Implicações dos slides que completam o quadro}

\begin{itemize}[leftmargin=1.4em]
\item \textbf{Poupança}: o padrão temporal da renda é irrelevante para o
consumo, mas \emph{crítico} para a poupança:
\[
S_t = Y_t - C_t
    = \Bigl(Y_t - \frac{1}{T}\sum_{\tau=1}^{T} Y_\tau\Bigr) - \frac{1}{T}A_0.
\]
Poupa-se quando a renda corrente está acima da permanente (o consumidor usa
poupança e empréstimo para suavizar o consumo).
\item \textbf{Regressões consumo--renda}: se $C = Y^P$ e $Y = Y^P + Y^T$, o
coeficiente estimado de uma regressão de $C$ sobre $Y$ é
\[
\hat{b} = \frac{\VAR(Y^P)}{\VAR(Y^P) + \VAR(Y^T)}.
\]
Amostras \textbf{longas}: variância permanente domina, $\hat{b}$ próximo de 1.
Amostras \textbf{curtas}: variância transitória pesa mais, $\hat{b}$ bem
abaixo de 1. Muito do desacordo Keynes--Friedman reflete apenas
\emph{horizontes de tempo diferentes}.
\item \textbf{Brasil} (slides): consumo das famílias e PIB caminham juntos no
longo prazo --- correlação de $0{,}99$ entre os índices reais (SCN/IBGE,
1996T1--2025T1).
\end{itemize}

\begin{provabox}[title={Como cai na prova}]
P2 2021 Q1-i / P2 2024 Q1-i (\textbf{falsos}): itens que trocam qual variância
domina. Guarde: \emph{amostras longas $\Rightarrow$ $\VAR(Y^P)$ domina
$\Rightarrow$ propensão estimada perto de 1}; a propensão \emph{baixa} aparece
no curto prazo, quando domina a variância \emph{transitória}. E o caput da
P2 2024 Q1-i está certo em chamar a teoria keynesiana de \emph{caso
particular} da renda permanente --- o erro do item estava na segunda parte.
\end{provabox}

% ------------------------------------------------------------------
\section{Passeio Aleatório (Hall, 1978) e o Teste de Campbell e Mankiw (1989)}
% ------------------------------------------------------------------

\subsection{Lista 3, Q1 (itens b e c) --- incerteza e teste empírico}

Com incerteza sobre a renda, utilidade \textbf{quadrática}
$\E[U] = \E\bigl[\sum_{t}\bigl(C_t - \tfrac{a}{2}C_t^2\bigr)\bigr]$, $a>0$, e
juros e desconto nulos, a solução do problema é
\[
C_1 = \frac{1}{T}\Bigl(A_0 + \sum_{t=1}^{T}\E_1[Y_t]\Bigr)
\]
--- a renda permanente, agora \emph{esperada}. Sob expectativas racionais,
segue a condição de otimização $C_t = \E_t[C_{t+1}]$.

\begin{correctbox}[title={Chave oficial (Q1b): a lógica de Hall}]
\emph{``A teoria do passeio aleatório é consistente com a teoria da renda
permanente. Hall (1978) parte da ideia de que as variações futuras do consumo
não podem ser previstas a partir de dados correntes, já que estes já foram
utilizados para a previsão da renda permanente pelos agentes.''}
\[
C_t = \E_t[C_{t+1}]
\quad\Longrightarrow\quad
C_t = \E_{t-1}[C_t] + e_t
\quad\Longrightarrow\quad
\boxed{C_t = C_{t-1} + e_t}
\]
\emph{``As variações entre $C_t$ e $C_{t-1}$ só podem ser observadas ex post e
são determinadas por informação nova em $t$ (a respeito da renda permanente).
Logo, não haveria forma de prever ex ante, em $t-1$, o nível de consumo efetivo
em $t$.''}

Atenção ao fundamento: o consumo é imprevisível \textbf{porque toda a
informação disponível já está incorporada} ao consumo corrente --- não porque
os agentes ignorem informação.
\end{correctbox}

\begin{correctbox}[title={Chave oficial (Q1c): o teste de Campbell e Mankiw}]
Regressão testada, com fração $\lambda$ de famílias ``regra de bolso'':
\[
C_t - C_{t-1} = \lambda\,(Y_t - Y_{t-1}) + (1-\lambda)\,e_t.
\]
\emph{``Campbell \& Mankiw (1989) testaram a hipótese do passeio aleatório de
Hall por meio de \textbf{variáveis instrumentais} (IV), desta maneira
controlando o problema de \textbf{endogeneidade} dos regressores, em particular
na variável explicativa $(Y_t - Y_{t-1})$.''}

Resultado (dados trimestrais dos EUA, 1953--1986): $\lambda$
\textbf{estatisticamente significativo, porém bem abaixo de 1}
($\approx 0{,}5$). Interpretação da chave: \emph{``explicação híbrida do
consumo''} --- uma fração $\lambda$ das famílias reage de modo estável à renda
corrente (à la Keynes), enquanto a parcela $(1-\lambda)$ comporta-se conforme
a renda permanente e o passeio aleatório.
\end{correctbox}

\begin{provabox}[title={Como cai na prova}]
Item \textbf{correto} da P2 2021 Q1 e da P2 2024 Q1 (gabaritos oficiais):
\emph{``a rejeição plena da hipótese do passeio aleatório requer encontrar
valores convergentes com $\lambda = 1$''} --- \textbf{verdadeiro}, pois
$\lambda = 1$ significaria que \emph{todas} as famílias seguem a renda
corrente. Como acharam $\lambda$ significativo mas $\ll 1$, a rejeição foi
apenas \emph{parcial}. Distrator recorrente (\textbf{falso}): ``a correlação
entre juros reais e crescimento do consumo implica o \emph{abandono} da
hipótese da renda permanente'' --- não implica: a equação de Euler é uma
\emph{extensão} do mesmo arcabouço otimizador.
\end{provabox}

% ------------------------------------------------------------------
\section{Juros Reais e Crescimento do Consumo (Equação de Euler)}
% ------------------------------------------------------------------

\subsection{Lista 3, Q2 (com chave oficial) --- teoria positiva, evidência fraca}

Com taxa de juros não nula, a restrição é em valor presente e a solução do
problema de otimização gera a equação de Euler:

\begin{formulabox}[title={Equação de Euler (forma da chave)}]
\[
\frac{C_{t+1}}{C_t} = \left(\frac{1+r}{1+\rho}\right)^{\frac{1}{\theta}}
\]
$r$ = taxa real de juros; $\rho$ = taxa de desconto subjetiva; $\theta$ = grau
de aversão relativa ao risco; $1/\theta$ = elasticidade de substituição
intertemporal (EIS).\\[0.3em]
$r > \rho$ $\Rightarrow$ consumo \emph{crescente} no tempo;\quad
$r < \rho$ $\Rightarrow$ consumo \emph{decrescente};\quad
sensibilidade do crescimento do consumo a $r$ dada por $1/\theta$.
\end{formulabox}

\begin{correctbox}[title={Chave oficial (Q2): como a literatura explica a fragilidade empírica}]
\emph{``Coeteris paribus, há uma relação teórica positiva entre crescimento do
consumo e taxa real de juros. Os autores mencionados}
[Hansen e Singleton, 1983; Hall, 1988b; Campbell e Mankiw, 1989]
\emph{associam o fato empírico de baixo poder preditivo de juros reais sobre
$C_{t+1}/C_t$ ao suposto grau elevado de aversão ao risco ($\theta$).''}

Em outras palavras: $\theta$ alto $\Rightarrow$ EIS $=1/\theta$ baixa
$\Rightarrow$ mesmo grandes variações de $r$ movem pouco a trajetória do
consumo. A \emph{teoria} não é abandonada; o \emph{parâmetro} é que torna o
canal fraco.
\end{correctbox}

\begin{provabox}[title={Como cai na prova}]
P2 2024 Q2: itens \textbf{falsos} típicos --- (i) ``variação do consumo
inversamente relacionada a $r$'' (é \emph{positiva}); (ii) ``há grande impacto
dos juros reais no consumo'' (o fato estilizado é impacto \emph{fraco});
(iii) ``maior $\rho$ e maior $\theta$ elevam a eficácia dos juros'' (maior
$\theta$ \emph{reduz} a sensibilidade $1/\theta$). O gabarito oficial daquela
questão (item iv) era a restrição orçamentária com governo --- ver Tópico 14
(equivalência ricardiana).
\end{provabox}

% ------------------------------------------------------------------
\section{Poupança Precaucional, o Puzzle e as Restrições de Liquidez}
% ------------------------------------------------------------------

\subsection{Lista 3, Q3 (com chave oficial) --- incerteza e poupança}

\begin{correctbox}[title={Chave oficial (Q3): motivo precaucional e conciliação com os fatos}]
\emph{``A poupança precaucional é explicada pela incerteza a respeito do
crescimento do consumo pelas famílias. (...) Quanto mais elevada a
variabilidade (especificamente do crescimento do consumo --- $\VAR(g^c)$),
maior a incerteza quanto ao futuro caminho da utilidade das famílias. E para
se defenderem disto há um incentivo à poupança pelo motivo de precaução. Desta
maneira, havendo esta poupança adicional, eleva-se a expectativa de
crescimento do consumo'':}
\[
\E[g^c] \simeq \tfrac{1}{2}\,(\theta + 1)\,\VAR(g^c).
\]
\emph{``Observe que há uma diferença entre o crescimento do consumo esperado e
sua variabilidade. A previsão é de que esta última impacta no primeiro.''}

\textbf{O puzzle}: \emph{``o fato estilizado é de que as famílias poupam
relativamente pouco, ou abaixo do que o modelo pressupõe. Para Carroll (1992,
1997), isto seria explicado pela elevada taxa de desconto subjetiva, que
acabaria anulando o aumento da poupança pelo motivo precaucional.''}
\end{correctbox}

Nos slides, a mesma expressão aparece como
$\E[g^c] \simeq \tfrac{1}{2}\,\delta\,\VAR(g^c)$, em que $\delta$ é o
\textbf{coeficiente de prudência relativa} --- para utilidade CRRA,
$\delta = \theta + 1$. Prudência ($u''' > 0$) é o que gera o motivo
precaucional; não confundir com aversão ao risco ($\theta$).

\subsection{Restrições de liquidez (slides da Aula 10)}

A HRP supõe que a família toma emprestado à mesma taxa a que poupa. Na
prática, as taxas de captação das famílias (cartão, financiamentos) são bem
maiores, e há limites de crédito. Restrições de liquidez elevam a poupança por
duas vias:
\begin{enumerate}[leftmargin=1.6em]
\item quando a restrição está \textbf{ativa} (\emph{binding}), a família
consome menos do que gostaria;
\item \textbf{Zeldes (1989)}: mesmo restrições \emph{não} ativas hoje reduzem
o consumo corrente, se puderem vir a ser ativas no futuro (evidência em Gross
e Souleles, 2002).
\end{enumerate}
\begin{center}
$\uparrow$ restrição de liquidez (efetiva \emph{ou esperada}) $\Rightarrow$
$\downarrow$ consumo.
\end{center}

\begin{provabox}[title={Como cai na prova}]
Distrator já usado \emph{duas vezes} (P2 2021 Q1-iv = P2 2024 Q1-iv,
\textbf{falso}): a expressão da poupança precaucional \emph{``ajuda a entender
o fato estilizado de uma elevada poupança pelas famílias''}. O fato estilizado
é o \textbf{oposto} --- poupança \emph{baixa} (por isso é um puzzle, resolvido
por Carroll via alta taxa de desconto). Se aparecer ``elevada poupança'',
marque falso sem hesitar.
\end{provabox}

% ==========================================================================
% PARTE II — INVESTIMENTO
% ==========================================================================
\newpage
\begin{center}
{\LARGE\bfseries\color{orangerbc} Parte II --- Teoria do Investimento e $q$ de Tobin}\\[0.2em]
{\normalsize\color{grayrbc} Aula 11 $\cdot$ Romer cap.\ 9 $\cdot$ Lista 3: Q4 e Q5}\\[0.2em]
\rule{\linewidth}{0.8pt}
\end{center}

% ------------------------------------------------------------------
\section{Custos de Ajuste e a Condição de Otimalidade da Firma}
% ------------------------------------------------------------------

\subsection{Por que estudar investimento (slides)}

A demanda por investimento importa para o \emph{longo prazo} (junto com a
poupança, determina quanto do produto é investido --- padrão de vida) e para o
\emph{curto prazo} (é o componente mais \textbf{volátil} da demanda --- nos
dados brasileiros dos slides, a FBCF oscila muito mais que o consumo).

\subsection{O modelo com custos de ajuste}

Hipótese-chave: a firma incorre em custos de ajuste \textbf{convexos} do
estoque de capital, $C(k)$ com $C(0)=0$, $C'(0)=0$ e $C''(\cdot)>0$. Os lucros
em cada instante são $\pi(K(t))\,k(t) - I(t) - C(I(t))$ --- lucros unitários
$\pi(\cdot)$ dependem do capital \emph{da indústria} $K$, com $\pi'(K)<0$
(mais capital agregado $\Rightarrow$ menor preço relativo do produto
$\Rightarrow$ menores lucros marginais). A firma maximiza o valor presente
\[
\Pi = \int_{t=0}^{\infty}
e^{-rt}\bigl[\pi(K(t))\,k(t) - I(t) - C(I(t))\bigr]\,\dif t.
\]

\begin{formulabox}[title={Condição de otimalidade e investimento agregado}]
Sendo $q_t$ o valor marginal (nominal) de uma unidade adicional de capital:
\[
1 + C'(I_t) = q_t
\qquad\Longrightarrow\qquad
\dot{K}(t) = f\bigl(q(t)\bigr),\ \ f(1) = 0,\ f'(\cdot)>0.
\]
Leitura dos slides: \emph{``a firma investe até o ponto em que o custo de
aquisição do capital (preço de compra --- fixado em 1 --- mais o custo
marginal de ajuste) \textbf{iguala} o valor do capital ($q_t$)''}.\\[0.3em]
$q>1$ $\Rightarrow$ investir ($\dot K > 0$);\quad
$q<1$ $\Rightarrow$ desinvestir ($\dot K < 0$);\quad
$q=1$ $\Rightarrow$ $\dot K = 0$.
\end{formulabox}

\begin{provabox}[title={Como cai na prova}]
Distrator repetido (P2 2021 Q2-i = P2 2024 Q3-i, \textbf{falso}): \emph{``a
firma investe até o ponto em que o valor do capital \underline{supera} o custo
de sua aquisição''}. No ótimo há \textbf{igualdade} --- se o valor superasse o
custo, haveria lucro marginal não explorado. A palavra ``supera'' é o erro.
\end{provabox}

% ------------------------------------------------------------------
\section{O $q$ de Tobin: Determinação, Produto e Juros}
% ------------------------------------------------------------------

\subsection{Lista 3, Q4 (com chave oficial) --- o que determina o $q$}

\begin{correctbox}[title={Chave oficial (Q4): definição e os dois canais}]
\emph{``O q de Tobin pode ser representado pela expressão''}
\[
q(t) = \int_{t=0}^{\infty} e^{-rt}\,\pi(K(t))\,\dif t,
\]
\emph{``sendo $q(t)$ o valor presente dos fluxos futuros de lucros da firma,
estes sujeitos à evolução do estoque de capital da indústria $K$.''}

\textbf{(a) Variações esperadas do produto}: \emph{``entram no cálculo do valor
de mercado das empresas por meio de seus impactos nos fluxos esperados de
lucros e de caixa. Dada a estrutura de juros, quanto maior o crescimento do
PIB, coeteris paribus, maior tende a ser o q, e logo maior tende a ser o valor
de mercado das empresas de um país''} --- melhora do retorno esperado e da
atratividade para investimentos.

\textbf{(b) Variações esperadas dos juros}: \emph{``possuem relação inversa
para com o q. A taxa real de juros (em particular os juros reais \textbf{de
longo prazo}) é a taxa de desconto, a taxa que traz os fluxos de caixa
esperados a valor presente. Logo, uma elevação da estrutura de juros na
economia, coeteris paribus, reduz o valor de mercado das empresas.''}
\end{correctbox}

\begin{formulabox}[title={Relação de consistência (locus $\dot q = 0$)}]
\[
\dot{q}(t) = r\,q(t) - \pi\bigl(K(t)\bigr)
\qquad\Longrightarrow\qquad
\dot q = 0 \iff q = \frac{\pi(K)}{r}.
\]
$\dot q$ é a diferença entre o custo de oportunidade do capital ($rq$) e seu
retorno corrente ($\pi(K)$). Como $\pi'(K)<0$: o \emph{locus} $\dot q = 0$ é
\textbf{negativamente inclinado} no plano $q \times K$; e uma queda de $r$
eleva $q$ para cada $K$ e torna a função $q=(1/r)\pi(K)$ \emph{mais
inclinada}.
\end{formulabox}

\begin{provabox}[title={Como cai na prova}]
P2 2024 Q4 (gabarito oficial: \textbf{iii}): \emph{``coeteris paribus, mais
importante do que os juros reais no momento $t$ é a \textbf{expectativa de $r$
para o longo prazo}, no que respeita ao q da empresa''} --- verdadeiro, direto
da chave (o desconto é a curva de juros toda, não a taxa de hoje). Falsos
típicos da mesma questão: valor da empresa \emph{positivamente} relacionado a
$K$ (é negativo, $\pi'<0$); \emph{locus} $\dot q=0$ \emph{positivamente}
inclinado (é negativo); e ``inflação esperada acima da meta $\Rightarrow$
maior $q$'' (falso: com regra de juros ativa, inflação acima da meta
$\Rightarrow$ juros reais maiores $\Rightarrow$ $q$ \emph{menor}).
\end{provabox}

% ------------------------------------------------------------------
\section{Diagrama de Fases, Trajetória de Sela e Choques}
% ------------------------------------------------------------------

\subsection{A dinâmica conjunta de $K$ e $q$}

\begin{resultbox}[title={Os dois loci e as setas de movimento}]
\textbf{Locus $\dot K = 0$}: reta horizontal em $q = 1$.
Acima ($q>1$): $K$ cresce ($\rightarrow$); abaixo ($q<1$): $K$ cai
($\leftarrow$).\\[0.3em]
\textbf{Locus $\dot q = 0$}: curva $q = \pi(K)/r$, decrescente.
À \emph{direita} (K grande, $\pi$ pequeno, $rq > \pi(K)$): $\dot q > 0$
($\uparrow$); à \emph{esquerda}: $\dot q < 0$ ($\downarrow$).\\[0.3em]
O equilíbrio $E$ (interseção) é um \textbf{ponto de sela}: dado $K$ inicial,
só há um valor de $q$ --- sobre a \textbf{trajetória de sela}, negativamente
inclinada --- consistente com a maximização de lucros; qualquer outro $q$
coloca a indústria em trajetória explosiva.
\end{resultbox}

\begin{correctbox}[title={Chave oficial (Q5): a situação ``acima de $\dot K=0$ e à esquerda de $\dot q=0$''}]
\emph{``Na situação proposta, a indústria está na \textbf{região de trajetória
de sela}.''}

\emph{``Em relação a $K$: seja a equação $\dot K(t) = NC'^{-1}(q-1)$. Uma vez
que acima de $\dot K = 0$, $q > 1$, então $\dot K(t) > 0$.''} (capital
crescendo)

\emph{``Em relação a $q$: seja $q = \pi(K)/r$. Uma vez que à esquerda de
$\dot q = 0$, o $q$ é menor do que o sugerido pelos lucros correntes, há uma
situação de \textbf{pessimismo quanto ao futuro}, e assim $\dot q < 0$.''}
($q$ caindo)

\emph{``A resultante dessas forças cria um vetor que leva a indústria rumo ao
ponto E''} --- ou seja, \textbf{há convergência} ao equilíbrio de longo prazo
pela trajetória de sela.
\end{correctbox}

\subsection{Choques (Figuras 9.5--9.7 dos slides)}

\begin{itemize}[leftmargin=1.4em]
\item \textbf{Aumento permanente do produto} (Fig.\ 9.5): demanda maior
$\Rightarrow$ lucros maiores para cada $K$ $\Rightarrow$ o \emph{locus}
$\dot q = 0$ desloca-se para fora. Como $K$ não se ajusta instantaneamente,
$q$ \textbf{salta} para a nova trajetória de sela; o investimento sobe; com a
acumulação, o preço relativo do produto e os lucros declinam até o valor do
capital voltar ao normal ($q = 1$) no novo $E'$ com $K$ maior.
\item \textbf{Aumento transitório} (Fig.\ 9.6): investidores sabem que o lucro
extra dura pouco $\Rightarrow$ o salto de $q$ e o efeito sobre o investimento
são \textbf{menores} do que no choque permanente --- mas \textbf{não nulos}.
(A indústria vai a um ponto \emph{abaixo} da nova sela e retorna à antiga.)
\item \textbf{Queda permanente dos juros} (Fig.\ 9.7): $q = (1/r)\pi(K)$ sobe
e fica mais inclinada; $q$ salta ($>1$), $\dot K > 0$ temporariamente, e a
indústria converge ao novo $E'$ com $K$ maior. O $r$ relevante é de
\textbf{longo prazo}.
\item \textbf{Acelerador}: \emph{``o investimento é maior quando o produto
subiu recentemente do que quando está alto há um período prolongado''} --- é a
\emph{variação} do produto que move o investimento.
\end{itemize}

\begin{provabox}[title={Como cai na prova}]
Item \textbf{verdadeiro} nas duas provas (P2 2021 Q2-iv = P2 2024 Q3-iv,
gabarito oficial): \emph{``pela equação de variação do q de Tobin, uma
situação de estoque de capital da indústria muito baixo, acompanhada de nível
também muito reduzido para o q, implica uma \textbf{queda adicional} do
último''} --- confira: $\dot q = rq - \pi(K)$ com $q$ baixo ($rq$ pequeno) e
$K$ baixo ($\pi(K)$ grande) dá $\dot q < 0$. É a região divergente
\emph{abaixo} da trajetória de sela.\\
Falsos recorrentes: ``região abaixo de $\dot K=0$ e à direita de $\dot q=0$ é
de divergência crescente'' (contém o braço convergente da sela); ``choque
transitório tem efeito \emph{nulo}'' (é \emph{menor}, nunca nulo).
\end{provabox}

% ==========================================================================
% PARTE III — INFLAÇÃO E POLÍTICA MONETÁRIA
% ==========================================================================
\newpage
\begin{center}
{\LARGE\bfseries\color{orangerbc} Parte III --- Inflação e Política Monetária}\\[0.2em]
{\normalsize\color{grayrbc} Aula 12 $\cdot$ Romer cap.\ 11 $\cdot$ Lista 3: Q6}\\[0.2em]
\rule{\linewidth}{0.8pt}
\end{center}

% ------------------------------------------------------------------
\section{Mercado Monetário, Nível de Preços e Identidade de Fisher}
% ------------------------------------------------------------------

\begin{formulabox}[title={Equações-base (slides, eqs.\ 1--4)}]
\[
\frac{M}{P} = L(i, Y),\quad L_i<0,\ L_Y>0
\qquad\Longrightarrow\qquad
P = \frac{M}{L(i,Y)}
\]
\[
i = r + \pi^e \ \text{(identidade de Fisher)}
\qquad\Longrightarrow\qquad
P = \frac{M}{L(\bar r + \pi^e,\, \bar Y)} \ \text{(preços flexíveis)}
\]
\end{formulabox}

Leituras:
\begin{itemize}[leftmargin=1.4em]
\item O nível de preços sobe com: $M\uparrow$; $i\uparrow$ (demanda por moeda
cai); $Y\downarrow$ (demanda por moeda cai). \textbf{O produto entra
negativamente} --- distrator clássico.
\item Para a inflação de \emph{longo prazo}, a ênfase recai em um único fator:
o \textbf{crescimento da oferta monetária} --- nenhum outro fator gera
aumentos \emph{persistentes} do nível de preços. Evidência (Fig.\ 11.1):
inflação média $\times$ crescimento monetário médio, 97 países, 1980--2006 ---
relação clara e forte.
\item Com preços plenamente flexíveis, moeda não afeta $Y$ nem $r$
(\textbf{neutralidade}; $\bar Y$, $\bar r$ constantes).
\end{itemize}

% ------------------------------------------------------------------
\section{Mudança Permanente no Crescimento da Moeda}
% ------------------------------------------------------------------

\subsection{Caso 1: preços plenamente flexíveis (longo prazo)}

Aumento permanente do crescimento da moeda em $t_0$:
\[
\uparrow \Delta M \;\Rightarrow\; \uparrow \pi^e \;\Rightarrow\; \uparrow i
\;\Rightarrow\; \downarrow \frac{M}{P}.
\]
Como a demanda por moeda real cai \emph{descontinuamente} e $M$ não salta, o
\textbf{nível de preços $P$ salta para cima} no instante da mudança
(Fig.\ 11.2). Duas mensagens dos slides:
\begin{enumerate}[leftmargin=1.6em]
\item \textbf{Efeito Fisher}: a inflação extra reflete-se um-para-um na taxa
nominal (a real não muda com preços flexíveis);
\item Crescimento maior da moeda \emph{nominal} $\Rightarrow$ estoque
\emph{real} de moeda \textbf{menor} no longo prazo ($P$ passa a crescer mais
que $M$ na transição).
\end{enumerate}

O espelho (redução permanente do crescimento da moeda) foi o item
\textbf{correto} da P2 2021 Q3-ii: queda descontínua de $P$, de $\pi^e$ e de
$i$, com \emph{aumento} de $M/P$ no longo prazo.

\subsection{Caso 2: flexibilidade incompleta de preços (curto prazo)}

\begin{formulabox}[title={Efeito liquidez vs.\ longo prazo (slides)}]
\textbf{Curto prazo:}\quad
$\uparrow \Delta M \Rightarrow \bar P \Rightarrow \uparrow \dfrac{M}{P}
\Rightarrow \downarrow r,\ \overline{\pi^e}
\Rightarrow \downarrow i \Rightarrow \uparrow Y$
\hfill (\emph{efeito liquidez})\\[0.4em]
\textbf{Longo prazo:}\quad
$\uparrow \Delta M \Rightarrow \uparrow \Delta P \Rightarrow
\uparrow \pi^e,\ \bar r \Rightarrow \uparrow i
\Rightarrow \downarrow \dfrac{M}{P}$
\end{formulabox}

Com preços rígidos no curto prazo, a expansão monetária eleva o estoque
\emph{real} de moeda; o reequilíbrio do mercado monetário exige queda da taxa
real (e nominal) de juros --- se o efeito taxa-real domina o efeito
inflação-esperada, a taxa nominal \textbf{cai no curto prazo} e \textbf{sobe no
longo prazo}. O benefício real é transitório; o efeito permanente é inflação
mais alta (exemplo citado em sala: Argentina).

% ------------------------------------------------------------------
\section{Estrutura a Termo das Taxas de Juros}
% ------------------------------------------------------------------

\begin{formulabox}[title={Sem e com incerteza (slides, eqs.\ 5--6)}]
Sem incerteza (arbitragem pura):
\[
i_t^n = \frac{i_t^1 + i_{t+1}^1 + \cdots + i_{t+n-1}^1}{n};
\]
com incerteza:
\[
i_t^n = \frac{i_t^1 + \E_t i_{t+1}^1 + \cdots + \E_t i_{t+n-1}^1}{n}
        + \theta_{nt},
\]
sendo $\theta_{nt}$ o \textbf{prêmio a termo} por carregar o título longo
(estratégias com riscos diferentes).
\end{formulabox}

\begin{itemize}[leftmargin=1.4em]
\item \textbf{Teoria das expectativas}: mudanças na estrutura a termo são
determinadas por mudanças nas \emph{expectativas das taxas curtas futuras}
(não no prêmio a termo); expectativas tipicamente racionais.
\item Mesmo com preços não plenamente flexíveis, um aumento permanente do
crescimento da moeda \emph{acaba} elevando a taxa curta nominal
permanentemente; logo, para $n$ suficientemente grande, as taxas
\textbf{longas tendem a subir imediatamente}.
\item Conclusão dos slides: \emph{``uma expansão monetária tende a
\textbf{reduzir} as taxas curtas} [efeito liquidez] \emph{mas \textbf{elevar}
as longas''}.
\end{itemize}

\begin{provabox}[title={Como cai na prova}]
Distrator da P2 2021 Q3-iii (\textbf{falso}): arbitragem exigiria
$i^n = n\,i^1$ --- errado, $i^n$ é \emph{média} das curtas esperadas mais o
prêmio. E o item iv da mesma prova combinava efeito liquidez com ``alta do
prêmio a termo'' --- a alta das longas vem das \emph{expectativas} de juros
curtos futuros maiores (inflação), não de um salto do prêmio.
\end{provabox}

% ------------------------------------------------------------------
\section{Custos da Inflação (Fundamentos Microeconômicos)}
% ------------------------------------------------------------------

Os seis custos listados nos slides (o objetivo do formulador de política é
maximizar bem-estar):
\begin{enumerate}[leftmargin=1.6em, itemsep=0.15em]
\item \textbf{Custo de sola de sapato}: inflação alta leva as pessoas a
gastarem esforço para reduzir encaixes de moeda de alto poder;
\item \textbf{Distorções de preços relativos} e má alocação de recursos;
\item \textbf{Distorção do sistema tributário} (Feldstein, 1997): a inflação
pesa mais nos preços de alimentos, com maior peso no orçamento das famílias de
baixa renda --- carga tributária mais \emph{regressiva};
\item \textbf{Perda de eficiência}: modelos formais (Bénabou, 1992; Tommasi,
1994);
\item \textbf{Dificuldade de contabilizar a inflação}: famílias e firmas erram
contas em ambiente inflacionário (Modigliani e Cohn, 1979; Hall, 1984);
\item \textbf{Relação média--variância}: inflação alta é também inflação
volátil; volatilidade (ou mesmo o nível alto) desestimula o investimento de
longo prazo, por sinalizar governo funcionando mal.
\end{enumerate}

% ------------------------------------------------------------------
\section{A Regra \emph{Forward-Looking} de Clarida, Galí e Gertler (2000)}
% ------------------------------------------------------------------

\subsection{Lista 3, Q6 (com chave oficial) --- a novidade de 2026}

Esta questão \textbf{não existia na Lista 3 de 2024}: é o acréscimo de 2026,
com chave detalhada --- candidata natural a questão inédita da P2.

\begin{correctbox}[title={Chave oficial (Q6-i): formulação e parâmetros}]
\[
i_t = r^n + \E_t(\pi_{t+n})
      + \phi_\pi\bigl[\E_t(\pi_{t+n}) - \pi^*\bigr]
      + \phi_y\,\E_t\bigl[\ln Y_{t+n} - \ln Y^n_{t+n}\bigr]
\]
\emph{``A regra estabelece que a taxa nominal de juros deve reagir às
expectativas de inflação e ao hiato esperado do produto.''}
\begin{itemize}[leftmargin=1.5em, itemsep=0.1em]
\item $i_t$: taxa nominal definida pela autoridade monetária;
\item $r^n$: taxa real \textbf{natural} (ou neutra) de juros;
\item $\E_t(\pi_{t+n})$: inflação esperada no horizonte relevante da política;
\item $\pi^*$: meta de inflação;
\item $\phi_\pi$: resposta aos desvios da inflação esperada frente à meta;
\item $\phi_y$: resposta ao hiato esperado do produto;
\item $\ln Y_{t+n} - \ln Y^n_{t+n}$: hiato esperado do produto.
\end{itemize}
É uma extensão da regra de Taylor (1993) com visão \emph{forward-looking} do
Banco Central. No equilíbrio de médio prazo ($\E_t[\pi_{t+n}]=\pi^*$, hiato
nulo): $i_t = r^n + \E_t[\pi_{t+n}]$.
\end{correctbox}

\begin{correctbox}[title={Chave oficial (Q6-ii): Princípio de Taylor e estabilidade da inflação}]
Fato estilizado: no longo prazo, economias com alto crescimento da oferta
monetária apresentam alta inflação média. A regra contribui para evitar esse
resultado ao estabelecer \emph{reação sistemática} a desvios da inflação. A
condição fundamental é o \textbf{Princípio de Taylor}:
\[
\phi_\pi > 1.
\]
\emph{``Essa condição implica que a taxa nominal aumenta mais do que
proporcionalmente a um aumento da inflação esperada''}; como
\[
\Delta r_t = \Delta i_t - \Delta \E_t(\pi_{t+n}),
\]
\emph{``quando $\phi_\pi > 1$, a taxa \textbf{real} de juros aumenta após um
choque inflacionário, reduzindo a demanda agregada e contribuindo para a
convergência da inflação em direção à meta''} --- política
\textbf{contracíclica} frente a movimentos da inflação esperada.

\emph{``Quando a autoridade monetária reage menos do que proporcionalmente
($\phi_\pi < 1$), a taxa real pode não se elevar o suficiente para conter as
pressões inflacionárias. Nessa situação, a inflação tende a tornar-se mais
\textbf{persistente} e as expectativas podem se \textbf{desancorar}.''}

Conclusão da chave: com $\phi_\pi > 1$, a política monetária ancora
expectativas, reduz a persistência inflacionária e garante inflação
\emph{baixa e estável} no longo prazo.
\end{correctbox}

\begin{resultbox}[title={Calibração citada em aula (linha do BCB, Relatórios Trimestrais)}]
$r^n = 5{,}0\%$;\quad $\pi^* = 3{,}0\%$;\quad $n = 4$ trimestres;\quad
$\phi_\pi = 1{,}5$;\quad $\phi_y = 0{,}3$:
\[
i_t = 5{,}0 + \E_t(\pi_{t+4})
      + 1{,}5\,[\E_t(\pi_{t+4}) - 3{,}0]
      + 0{,}3\,\E_t(\ln Y_{t+4} - \ln Y^n_{t+4}).
\]
Nas funções de resposta a impulso mostradas em aula, o cenário $\phi_\pi=1{,}5$
traz a inflação de volta \emph{mais rápido} que $\phi_\pi=0{,}7$ após um choque
inflacionário de 1 desvio-padrão --- ao custo de queda maior do produto no
curto prazo.
\end{resultbox}

\begin{provabox}[title={Como cai na prova}]
Prováveis pontos de julgamento V/F: (i) trocar \emph{forward} por
\emph{backward-looking} (a regra reage a \textbf{expectativas}, não à inflação
passada); (ii) princípio de Taylor com desigualdade invertida ou aplicado à
taxa \emph{nominal} (o que sobe mais que proporcionalmente é a nominal, para
que a \textbf{real} suba); (iii) equilíbrio de médio prazo:
$i = r^n + \pi^*$ (com a calibração, $8\%$) --- cuidado com sinais; (iv)
afirmar que $\phi_\pi<1$ só retarda a convergência (falso: gera persistência e
desancoragem).
\end{provabox}

% ==========================================================================
% PARTE IV — POLÍTICA FISCAL
% ==========================================================================
\newpage
\begin{center}
{\LARGE\bfseries\color{orangerbc} Parte IV --- Déficits Orçamentários e Política Fiscal}\\[0.2em]
{\normalsize\color{grayrbc} Aula 13 $\cdot$ Romer cap.\ 12 $\cdot$ Lista 3: Q7 e Q8}\\[0.2em]
\rule{\linewidth}{0.8pt}
\end{center}

% ------------------------------------------------------------------
\section{Restrição Orçamentária Intertemporal do Governo e Solvência}
% ------------------------------------------------------------------

\subsection{Fatos estilizados (slides)}

Vários países mantêm déficits elevados de forma \emph{sistemática}, com breves
períodos de superávit; o envelhecimento populacional (mais aposentados por
trabalhador) pressiona previdência e saúde; a percepção é de que déficits
grandes e persistentes reduzem o crescimento e podem levar a crises.
\textbf{Brasil} (slides, fonte BCB): resultado primário médio de
$-0{,}34\%$ do PIB e nominal médio de $-4{,}5\%$ do PIB entre 06/1999 e
05/2025; dívida bruta do governo geral saiu de $\approx 55\%$ (2006) para picos
próximos de $90\%$ do PIB (2020), rondando $75\%$ em 2025 --- pergunta do
professor em aula: \emph{``o Brasil tem operado sob um regime fiscal
sustentável?''}

\subsection{Lista 3, Q7 (com chave oficial)}

\begin{correctbox}[title={Chave oficial (Q7): a restrição e o papel das expectativas}]
\emph{``A restrição do governo pode ser entendida como''}
\[
(1)\quad \int_{t=0}^{\infty} e^{-R(t)}\,G(t)\,\dif t
\;\leq\; -D(0) + \int_{t=0}^{\infty} e^{-R(t)}\,T(t)\,\dif t,
\]
\emph{``estabelecendo a necessidade de que o somatório dos gastos públicos a
valor presente seja menor ou igual ao somatório de receitas com arrecadação, a
valor presente, menos a dívida pública inicial. Isto leva à \textbf{condição de
solvência} do setor público:''}
\[
(2)\quad \int_{t=0}^{\infty} e^{-R(t)}\bigl[T(t) - G(t)\bigr]\,\dif t
\;\geq\; D(0).
\]
\emph{``Em todo momento do tempo, a solvência é avaliada pela capacidade do
governo gerar fluxos corrente e futuros, a valor presente, suficientes para
cobrir os passivos e dívidas públicas. Logo, um aumento da dívida pública
inicial deve ser acompanhado de elevação do fluxo temporal esperado de
\textbf{superávits primários}.''}

\emph{``As \textbf{expectativas} aqui ganham relevância porque a identificação
de solvência depende muito mais dos fluxos \emph{projetados} de resultados
fiscais, em comparação com o resultado fiscal corrente (apenas uma pequena
fração do lado esquerdo na condição 2). Portanto, \textbf{anúncios de política
fiscal e sinalizações} por parte das autoridades fiscais têm especial
importância na avaliação de solvência pelos agentes.''}
\end{correctbox}

Implicações dos slides: (1) sustentabilidade fiscal depende não só dos
déficits correntes, mas das \emph{expectativas sobre a política fiscal
futura}; (2) papel essencial de \emph{expectativas fiscais, reputação e
credibilidade}. Há um caso para \emph{regras fiscais}: a restrição diz que o
governo precisa gerar superávits primários grandes o bastante, em valor
presente, para compensar a dívida inicial.

\begin{provabox}[title={Como cai na prova}]
P2 2024 Q6 (gabarito oficial: \textbf{iii}): \emph{``os superávits primários
corrente e futuros esperados precisam cobrir a dívida pública atual em termos
reais''} --- é a condição (2). Falsos da mesma questão: solvência avaliada
pelos \emph{resultados correntes} (i); país com pior saldo corrente
\emph{necessariamente} mais arriscado (ii --- expectativas podem inverter);
avaliação fiscal imune à dinâmica inflacionária (iv).
\end{provabox}

% ------------------------------------------------------------------
\section{Equivalência Ricardiana}
% ------------------------------------------------------------------

\subsection{Lista 3, Q8 (com chave oficial) --- derivação completa}

Ponto de partida: modelo \textbf{RCK com tributos lump-sum}. Com impostos, o
valor presente do consumo não pode exceder a riqueza inicial (capital +
títulos públicos) mais o valor presente da renda do trabalho \emph{após
impostos}.

\begin{correctbox}[title={Chave oficial (Q8): passo a passo}]
\textbf{(1) Restrição das famílias:}
\[
\int_{t=0}^{\infty} e^{-R(t)} C(t)\,\dif t
\;\leq\; K(0) + D(0)
+ \int_{t=0}^{\infty} e^{-R(t)}\bigl[W(t) - T(t)\bigr]\,\dif t
\]
(separando: $\dots + \int e^{-R}W\,\dif t - \int e^{-R}T\,\dif t$).

\textbf{(2) Restrição do governo com igualdade:}
\[
\int_{t=0}^{\infty} e^{-R(t)} G(t)\,\dif t
= -D(0) + \int_{t=0}^{\infty} e^{-R(t)} T(t)\,\dif t,
\]
que pode ser reescrita como
\[
(3)\quad \int_{t=0}^{\infty} e^{-R(t)} T(t)\,\dif t
= \int_{t=0}^{\infty} e^{-R(t)} G(t)\,\dif t + D(0).
\]

\textbf{Substituindo (3) em (1):}
\[
\boxed{\;\int_{t=0}^{\infty} e^{-R(t)} C(t)\,\dif t
\;\leq\; K(0)
+ \int_{t=0}^{\infty} e^{-R(t)} W(t)\,\dif t
- \int_{t=0}^{\infty} e^{-R(t)} G(t)\,\dif t\;}
\]
$D(0)$ e a trajetória de $T$ \textbf{desaparecem da restrição}.

\emph{``Portanto, dados os gastos públicos corrente e futuros, a forma de
financiamento dos mesmos (se via arrecadação T, ou via dívida D) não interfere
na determinação do consumo intertemporal; e dada a trajetória de longo prazo
do produto, nem nos investimentos, já que $I(t) = Y(t) - C(t) - G(t)$.
Portanto, substituir tributos por dívida hoje, ou vice-versa, \textbf{não
impacta a atividade econômica de curto prazo}.''}
\end{correctbox}

\begin{resultbox}[title={Implicações (slides) e o mecanismo econômico}]
\begin{enumerate}[leftmargin=1.6em, itemsep=0.15em]
\item A restrição das famílias é afetada pelo \textbf{valor presente das
compras do governo}, sem referência à divisão do financiamento entre tributos
e títulos em cada momento;
\item São as \textbf{compras do governo}, não os tributos, que afetam a
acumulação de capital ($I = Y - C - G$);
\item \textbf{Resultado-chave}: só a \emph{quantidade} de gasto público, não a
divisão do seu financiamento, afeta a economia.
\end{enumerate}
Mecanismo: um corte de tributos hoje, dados os gastos, eleva a dívida; agentes
racionais internalizam a restrição do governo e sabem que virão maiores
tributos/superávits no futuro --- a renda disponível extra (transitória, pela
Parte I deste guia!) é \textbf{poupada}, e o consumo não muda.
\end{resultbox}

\begin{provabox}[title={Como cai na prova}]
Já caiu de três formas: P2 2024 Q2-iv (\textbf{verdadeiro}): a restrição das
famílias \emph{``não sofre influência das proporções de T e D''}; P2 2024
Q5-iii (\textbf{verdadeiro}, gabarito oficial): satisfeita a restrição do
governo, \emph{``a substituição de financiamento via tributos por mais dívida
não possui qualquer impacto sobre o comportamento do consumidor''}; e como
dissertativa (P2 2021 e 2024). Falsos associados: ``redução de alíquotas
estimula consumo e ciclo'' (P2 2024 Q5-ii); ``gastos públicos não têm impacto
direto na restrição das famílias'' (Q5-iv --- têm: $G$ entra com sinal negativo
após a substituição).
\end{provabox}

% ------------------------------------------------------------------
\section{Contrações Fiscais Expansionistas e Crises de Dívida}
% ------------------------------------------------------------------

\subsection{Conteúdo novo dos slides (Aula 13) --- ainda não cobrado em prova}

\textbf{Contrações fiscais expansionistas} --- às vezes a consolidação fiscal
\emph{aumenta} a atividade econômica:
\begin{itemize}[leftmargin=1.4em]
\item \textbf{Giavazzi e Pagano (1990)}: pacotes de reforma fiscal na
\textbf{Dinamarca e na Irlanda nos anos 1980} foram seguidos de \emph{booms de
consumo}; os autores atribuem o resultado a efeitos operando via
\textbf{expectativas};
\item \textbf{Alesina e Perotti (1997)}: reduções de déficit vindas de
\textbf{cortes de emprego público e transferências} são muito mais prováveis
de serem \emph{mantidas} do que as vindas de aumentos de tributos; e,
coerentemente com o papel das expectativas, o primeiro tipo é com frequência
\emph{expansionista}, o segundo usualmente não;
\item os \textbf{multiplicadores fiscais} diferem entre os dois tipos de
ajuste.
\end{itemize}

\textbf{Canais} (slides):
\begin{enumerate}[leftmargin=1.6em, itemsep=0.1em]
\item[a)] \emph{Juros}: menos compras do governo reduzem juros correntes;
expectativa de menos compras futuras reduz as expectativas de juros futuros;
\item[b)] \emph{Lado da oferta}: menores tributos futuros $\Rightarrow$
menores distorções (maior competitividade das firmas) $\Rightarrow$ maior
renda futura;
\item[c)] \emph{Risco de crise}: a contração reduz a probabilidade percebida
de crise, elevando as estimativas de renda futura;
\item[d)] renda futura esperada maior eleva \textbf{consumo corrente e
investimento} (Bertola e Drazen, 1993; Perotti, 1999; via \emph{q de Tobin} ---
ponte com a Parte II).
\end{enumerate}

\textbf{Modelos de crises de dívida} (inspirados em Calvo, 1988; Cole e Kehoe,
2000) --- crises de dívida são \textbf{crises induzidas por expectativas}:
\begin{center}
dívida alta $\rightarrow$ taxa de retorno requerida alta $\rightarrow$\\
receitas futuras esperadas baixas $\rightarrow$ \emph{default} mais provável,
\end{center}
um círculo com potencial autorrealizável: o medo do default eleva os juros
exigidos, o que piora a dinâmica da dívida e aproxima o próprio default.

\begin{provabox}[title={Como cai na prova}]
Por ser conteúdo novo, itens plausíveis: inverter Alesina--Perotti (dizer que
ajustes via \emph{tributos} são os duradouros/expansionistas --- falso);
inverter o círculo de Calvo/Cole-Kehoe (dívida alta $\rightarrow$ retorno
requerido \emph{menor} --- falso); ou afirmar que consolidações são
\emph{sempre} contracionistas (falso: Dinamarca e Irlanda). A âncora
verdadeira: expectativas no centro dos dois fenômenos.
\end{provabox}

% ==========================================================================
% FECHAMENTO
% ==========================================================================
\newpage
\begin{center}
{\large\bfseries\color{orangerbc} Resumo de véspera --- as equações da P2 em meia página}\\[0.2em]
\rule{\linewidth}{0.6pt}
\end{center}

\begin{formulabox}[title={Consumo (cap.\ 8)}]
$C_t = \frac{1}{T}\bigl(A_0 + \sum_\tau Y_\tau\bigr)$
\quad$\cdot$\quad
$\hat b = \frac{\VAR(Y^P)}{\VAR(Y^P)+\VAR(Y^T)}$
\quad$\cdot$\quad
$C_t = C_{t-1} + e_t$
\quad$\cdot$\quad
$\Delta C_t = \lambda \Delta Y_t + (1-\lambda)e_t$ (IV; $\lambda\approx0{,}5$)
\quad$\cdot$\quad
$\frac{C_{t+1}}{C_t} = \bigl(\frac{1+r}{1+\rho}\bigr)^{1/\theta}$
\quad$\cdot$\quad
$\E[g^c] \simeq \frac{1}{2}(\theta+1)\VAR(g^c)$
\end{formulabox}

\begin{formulabox}[title={Investimento (cap.\ 9)}]
$q(t) = \int_0^\infty e^{-rt}\pi(K(t))\dif t$
\quad$\cdot$\quad
$1 + C'(I) = q$
\quad$\cdot$\quad
$\dot K = f(q),\ f(1)=0$
\quad$\cdot$\quad
$\dot q = rq - \pi(K)$
\quad$\cdot$\quad
$\dot q = 0 \Rightarrow q = \pi(K)/r$ ($\pi'<0$: inclinação negativa)
\end{formulabox}

\begin{formulabox}[title={Inflação e política monetária (cap.\ 11)}]
$P = M/L(i,Y)$
\quad$\cdot$\quad
$i = r + \pi^e$
\quad$\cdot$\quad
$i_t^n = \frac{1}{n}\sum_{j=0}^{n-1}\E_t\,i_{t+j}^1 + \theta_{nt}$
\quad$\cdot$\quad
$i_t = r^n + \E_t\pi_{t+n} + \phi_\pi(\E_t\pi_{t+n} - \pi^*) + \phi_y\E_t[\ln Y_{t+n} - \ln Y^n_{t+n}]$,
\ Princípio de Taylor: $\phi_\pi > 1$
\end{formulabox}

\begin{formulabox}[title={Política fiscal (cap.\ 12)}]
$\int_0^\infty e^{-R(t)}[T(t)-G(t)]\dif t \geq D(0)$
\quad$\cdot$\quad
$\int e^{-R}C \leq K(0) + \int e^{-R}W - \int e^{-R}G$ (equivalência
ricardiana: $T$, $D$ somem)
\quad$\cdot$\quad
$I = Y - C - G$
\end{formulabox}

\begin{provabox}[title={Roteiro final de estudo}]
\begin{enumerate}[leftmargin=1.6em, itemsep=0.15em]
\item Releia os 15 tópicos deste guia marcando as caixas verdes (chave
oficial) --- elas são o padrão de resposta do professor;
\item Faça o simulado (\texttt{provas/p2\_simulado\_2026.pdf}) cronometrado e
confira o gabarito comentado;
\item Na véspera, revise as caixas ``Como cai na prova'' e o resumo de
equações acima;
\item Na prova: leia as 4 afirmativas inteiras antes de marcar --- a diferença
entre V e F costuma ser \emph{uma palavra} (``supera''/``iguala'',
``nulo''/``menor'', ``corrente''/``esperado'', ``positiva''/``negativa'').
\end{enumerate}
\end{provabox}

\vfill
{\small\color{grayrbc} Guia de conteúdo da Prova 2 --- Macroeconomia I,
PPGEco-UFES 2026. Elaborado a partir da Lista 3 2026 e sua chave oficial, dos
slides das Aulas 10--13 e das provas P2 2021 e P2 2024 (chave oficial). Uso
pessoal.}

\end{document}
